\section{Topological and algebraic reductions}\label{sec:foundations}

Fix a finite connected aspherical integral Poincar\'e complex $Q$ of
formal dimension four, with the absolute duality convention of the
introduction, and put $\Gamma=\pi_1(Q)$ and $k=\mathbb F_2$.
This section reduces Theorem~\ref{thm:main} to a statement about functions
on a Cayley graph. We first justify the topological and coefficient
conventions, then extract the group-theoretic restrictions needed below.

\begin{lemma}\label{found:model}
The complex $Q$ is homotopy equivalent to a finite four-dimensional
CW complex. The group $\Gamma$ is infinite, finitely presented, and
torsion-free.

Every closed connected aspherical topological four-manifold is homotopy
equivalent to a finite connected aspherical integral Poincar\'e complex
of formal dimension four.
\end{lemma}

\begin{proof}
Wall's dimension theorem \cite[Theorem~2.2]{Wall1967} gives a
four-dimensional CW complex homotopy equivalent to $Q$. It can be
chosen finite because $Q$ is finite, so its reduced finiteness
obstruction vanishes. The model is aspherical, and its universal
cover is contractible. Its augmented cellular chain complex is
consequently a free $\mathbb Z\Gamma$-resolution of $\mathbb Z$ of
length four.
In particular, $\Gamma$ is finitely presented and has finite integral
cohomological dimension. To see directly why it has no torsion, a
finite-order element would yield a subgroup $C$ of prime order.
Restricting the resolution to $\mathbb Z C$ preserves freeness and
finite length. This contradicts the usual two-periodic resolution of
a cyclic group, which gives nonzero groups $H^{2j}(C;\mathbb Z)$ for
every $j\geq1$.
If $\Gamma$ were finite, torsion-freeness would make it trivial, so
asphericity would make $Q$ contractible. But duality over $k$ gives
$H_4(Q;k)\cong H^0(Q;k)\cong k$. Hence $\Gamma$ is infinite.

For the manifold assertion, let $M$ be a closed connected aspherical
topological four-manifold.
A compact topological manifold is metrizable. Its coordinate neighborhoods
are absolute neighborhood retracts, so Hanner's locality theorem makes
$M$ an absolute neighborhood retract; see
\cite[Theorems~3.2--3.3]{Hanner1951}.
To check its covering dimension, shrink a finite coordinate cover to a
finite closed cover. Each member has small inductive dimension at most
four, by the subspace theorem and the dimension of Euclidean space.
The closed-sum and coincidence theorems give $\dim M\le4$; conversely,
a coordinate chart contains a copy of an open four-ball, so subspace
monotonicity gives $\dim M\ge4$. These dimension facts are
\cite[Theorems~1.1.2, 1.5.3, 1.7.7 and~1.8.2]{Engelking1978}.
West's homotopy dimension estimate
\cite[Corollary~5.4]{West1977} therefore supplies a finite
four-dimensional simplicial complex homotopy equivalent to $M$.
This complex is a homotopy model, not a triangulation of the manifold.
It is aspherical. Transport the orientation character and fundamental
class of $M$ to this model. Topological Poincar\'e duality with the
orientation local system gives cap-product isomorphisms for every local
coefficient system on $M$ \cite[Theorem~1, with $R=\mathbb Z$]{Sun2017}.
Since $M$ is compact, compactly supported cohomology there is ordinary
cohomology. Transporting these isomorphisms along the homotopy equivalence
gives the required absolute integral duality on the model; see also
\cite[Theorem~10.2 and Remark~10.3]{Spanier1993}, and
\cite[Theorem~A.16]{FNOP2025} for the oriented chain version.
\end{proof}

\subsection{The singular-chain definition and duality}

We write $b_p^{(2)}(\widetilde Q)$ for the invariant defined using
$\mathcal N(\Gamma)\otimes_{\mathbb Z\Gamma}C_*^s(\widetilde Q)$.
The following comparison explains the passage from that definition to
finite cellular matrices and, later, to a presentation graph.

\begin{lemma}\label{found:comparison}
Let $\Lambda$ be a group and let $P_*$ be a free
$\mathbb Z\Lambda$-resolution of $\mathbb Z$ whose terms in degrees
$0,1,2$ are finitely generated. Suppose that a contractible space $Z$
has a free $\Lambda$-action for which its singular chains are free
$\mathbb Z\Lambda$-modules. Then
\[
 b_1^{(2)}(Z;\mathcal N(\Lambda))
 =\dim_{\mathcal N(\Lambda)}
   \bigl(\ker\partial_1\cap\ker\partial_2^*\bigr),
\]
where the operators on the right act on the finite Hilbert modules
$\ell^2(\Lambda)\otimes_{\mathbb Z\Lambda}P_i$.
The same space represents first reduced Hilbert cohomology for
$d_0=\partial_1^*$ and $d_1=\partial_2^*$.
The analogous comparison holds in every degree of a finite free
cellular complex.
\end{lemma}

\begin{proof}
The augmented singular chains of $Z$ are a free resolution of
$\mathbb Z$. The comparison theorem for free resolutions gives chain
homotopy equivalences with $P_*$.
Tensoring preserves these homotopies, so it preserves the homology in
the definition of $b_1^{(2)}$; no flatness of $\mathcal N(\Lambda)$
over the group ring is required.

The remaining assertion is the finite-matrix comparison for von
Neumann dimension. For a matrix over $\mathcal N(\Lambda)$, the
algebraic kernel has dimension equal to the trace of the projection
onto its Hilbert kernel. The algebraic image has dimension equal to
the trace of the projection onto the closure of its Hilbert image.
Additivity of dimension thus identifies the dimension of first
algebraic homology with that of
\[
 \ker\partial_1\ominus\overline{\operatorname{im}\partial_2}
 =\ker\partial_1\cap\ker\partial_2^*.
\]
On the other hand, the harmonic representative space for reduced
cohomology is
\[
 \ker d_1\ominus\overline{\operatorname{im}d_0}
 =\ker\partial_2^*\cap\ker\partial_1.
\]
These are the same closed Hilbert submodule. The same argument works
in each degree with finite free neighboring terms. This is the
comparison underlying \cite[Remark~1.31 and Lemmas~6.51--6.53]{Luck2002}.
\end{proof}

\begin{lemma}\label{found:reduction}
For the Poincar\'e complex $Q$ fixed above,
\[
 b_0^{(2)}(\widetilde Q)=b_4^{(2)}(\widetilde Q)=0,
 \qquad
 b_3^{(2)}(\widetilde Q)=b_1^{(2)}(\widetilde Q),
\]
and $b_p^{(2)}(\widetilde Q)=0$ for $p>4$.
Consequently it suffices to prove $b_1^{(2)}(\widetilde Q)=0$.
\end{lemma}

\begin{proof}
Use the finite four-dimensional model in Lemma~\ref{found:model}.
A homotopy equivalence with $Q$ lifts to an equivariant homotopy
equivalence of universal covers, and cellular and singular chains
compute the same invariant; see
\cite[Definition~6.50 and Lemmas~6.51--6.53]{Luck2002}.
There are no cellular chains above degree four, proving the assertion
for $p>4$. A harmonic square-summable $0$-cochain on the connected
universal cover is constant. Since $\Gamma$ is infinite, that constant
must be zero. The degree-zero version of
Lemma~\ref{found:comparison} gives $b_0^{(2)}=0$.

Here is the duality argument, including nontrivial orientation
characters. Let $\Gamma^+$ be the kernel of the orientation
character $\Gamma\to\{1,-1\}$, and let $Q^+$ be the corresponding
connected cover. Its degree $d=[\Gamma:\Gamma^+]$ is either one or
two. Finite-sheeted covers of Poincar\'e complexes are again
Poincar\'e complexes of the same formal dimension
\cite[p.~567, paragraph preceding Corollary~4.4]{Stark1996}. The induced orientation
character of $Q^+$ is the restriction to $\Gamma^+$ and is trivial.
Thus $Q^+$ is a finite connected aspherical oriented integral
Poincar\'e complex of formal dimension four, and
Lemma~\ref{found:model} supplies a finite four-dimensional homotopy
model $F$.

Transport the fundamental class of $Q^+$ to $F$.
Poincar\'e duality with local coefficients applies to $Q^+$ and hence
to $F$; see \cite[Lemma~1.1]{Wall1967}.
In particular, applying it to the regular group-ring local system
shows that the cap-product map from the degree-reversed dual of
$C_*(\widetilde F;\mathbb Z)$ to $C_*(\widetilde F;\mathbb Z)$
is a quasi-isomorphism. Here right modules are converted to left
modules by the involution $g\mapsto g^{-1}$.
Both complexes are bounded and finitely generated free over
$\mathbb Z\Gamma^+$. The mapping cone is therefore a bounded
acyclic complex of free modules, which is contractible: starting in
its lowest degree, split each surjection onto the kernel already
obtained. Thus cap product is a chain homotopy equivalence.

Extend scalars to $\mathcal N(\Gamma^+)$.
The homotopies are preserved, and finite-matrix dimension comparison
identifies dual Hilbert complexes with equal von Neumann dimensions.
It follows that
\[
 b_p^{(2)}(\widetilde Q;\mathcal N(\Gamma^+))
 =b_{4-p}^{(2)}(\widetilde Q;\mathcal N(\Gamma^+)),
 \qquad 0\leq p\leq4.
\]
This is the finite-chain argument underlying the $L^2$-duality proof
in \cite[Theorem~1.35(3) and its proof]{Luck2002}; here
Poincar\'e-complex duality supplies the required chain equivalence.
Finally, finite-index scaling
\cite[Theorem~1.35(9)]{Luck2002} gives
\[
 b_p^{(2)}(\widetilde Q;\mathcal N(\Gamma^+))
 =d\,b_p^{(2)}(\widetilde Q;\mathcal N(\Gamma)).
\]
Dividing by $d$ proves the same duality for $Q$ and completes the
reduction.
\end{proof}

\subsection{Cohomology, ends, and splittings}

For the boundary argument we use mod-two coefficients.
If $J\leq\Gamma$, write $k[\Gamma/J]$ for the left permutation
module on left cosets. Its elements are finite sums of cosets.

\begin{lemma}\label{found:duality}
The group $\Gamma$ satisfies
\[
 H^i(\Gamma;k\Gamma)=0\quad(i=1,2).
\]
For every subgroup $J\leq\Gamma$ there is an isomorphism
\[
 H^1(\Gamma;k[\Gamma/J])\cong H_3(J;k).
\]
In particular, this group vanishes when $J$ is finite or two-ended.
\end{lemma}

\begin{proof}
The orientation twist of $Q$ is trivial over $k$. Since $Q$ is
aspherical, its absolute integral Poincar\'e duality, applied to
local systems of $k$-vector spaces, gives, for
every left $k\Gamma$-module $V$,
\[
 H^i(\Gamma;V)\cong H_{4-i}(\Gamma;V).
\]
For clarity, the induction used on the right is homological.
Indeed, $k[\Gamma/J]=k\Gamma\otimes_{kJ}k$, and a free right
$k\Gamma$-resolution $R_*$ of $k$ satisfies
\[
 R_*\otimes_{k\Gamma}(k\Gamma\otimes_{kJ}k)
 \cong (R_*|_J)\otimes_{kJ}k.
\]
Restriction to $kJ$ preserves a free resolution. This proves
$H_j(\Gamma;k[\Gamma/J])\cong H_j(J;k)$, the homological
Shapiro isomorphism; compare \cite[Chapter~III, Section~6]{Brown1982}.
Taking $J$ trivial gives the assertion for $k\Gamma$, and taking
$i=1$ gives the displayed isomorphism for arbitrary $J$.

A finite subgroup is trivial by Lemma~\ref{found:model}.
A two-ended group has a finite normal subgroup with quotient
infinite cyclic or infinite dihedral; see \cite[Lemma~4.1]{Wall1967}.
If it is torsion-free the
finite subgroup is trivial and the dihedral case is impossible;
thus a two-ended subgroup of $\Gamma$ is infinite cyclic.
Both the trivial group and the infinite cyclic group have zero
third homology over $k$.
\end{proof}

\begin{lemma}\label{found:ends}
The group $\Gamma$ is one-ended.
\end{lemma}

\begin{proof}
Choose a finite presentation, and consider its simply connected
Cayley two-complex. Because there are finitely many cell orbits in
degrees at most two, finite-support cellular cochains in these
degrees identify with group-ring-valued equivariant cochains.
Explicitly, a finite-support $k$-cochain $a$ corresponds to
\[
 c\longmapsto\sum_{g\in\Gamma}a(g^{-1}c)g.
\]
Taking the identity coefficient recovers $a$, and the identification
commutes with coboundaries.
The augmented cellular complex is exact in degrees zero and one;
adjoining free modules in higher degrees extends it to a resolution
without changing its cohomology calculation in degree one.

Suppose that removing finitely many vertices of the Cayley graph
leaves two infinite components. Let $f$ be the indicator of one
infinite component, with value zero on the deleted vertices.
Its gradient $df$ is a finitely supported edge cochain, because
the graph is locally finite, and it is a cellular cocycle.
By Lemma~\ref{found:duality}, $df=du$ for a finitely supported
vertex function $u$. Connectedness makes $f-u$ constant.
This is impossible: $f$ takes each of its values on infinitely
many vertices, whereas $u$ has finite support.
The infinite connected locally finite Cayley graph has at least
one end, so it has exactly one.
\end{proof}

We also need a form of the splitting obstruction that applies to
tree actions arising during the proof. An action on a simplicial
tree is \emph{without inversions} if no element exchanges the
endpoints of an edge it preserves. It is \emph{minimal} if the
tree has no proper nonempty invariant subtree.

\begin{lemma}\label{found:trees}
For an action of $\Gamma$ without inversions on a simplicial tree:
\begin{enumerate}
\item if the action has no fixed vertex, its edge stabilizers
cannot all be finite or two-ended;
\item if the action is minimal and the tree has an edge, every
edge stabilizer is neither finite nor two-ended.
\end{enumerate}
Consequently $\Gamma$ has no nontrivial splitting over a finite
or two-ended subgroup.
\end{lemma}

\begin{proof}
Fix a vertex $o$ and an edge orbit $E$ whose stabilizer is $J$.
Identify $E$ with $\Gamma/J$, and define $c(g)$ to be the mod-two
sum of the edges of $E$ on the geodesic $[o,go]$.
Cancellation of edges traversed twice gives
\[
 c(gh)=c(g)+g c(h).
\]
If $J$ is finite or two-ended, Lemma~\ref{found:duality} makes
this cocycle principal: $c(g)=gv-v$ for some
$v\in k[\Gamma/J]$. Therefore
\begin{equation}\label{found:crossing-bound}
 |[o,go]\cap E|\leq2|\operatorname{supp}v|
 \qquad(g\in\Gamma).
\end{equation}

For the first assertion, let $S$ be a finite generating set of
$\Gamma$. The paths $[o,so]$, $s\in S$, contain finitely many
edges. Their translates contain every path $[o,go]$, by writing
$g$ as a word in $S$ and cancelling backtracking.
Thus only finitely many edge orbits occur in these paths.
An action without a fixed vertex has unbounded vertex orbits:
a bounded tree orbit has an invariant center, which gives a
fixed vertex when inversions are absent.
Hence one of the finitely many edge-orbit crossing counts is
unbounded. If all edge stabilizers were finite or two-ended,
this would contradict \eqref{found:crossing-bound}.

For the second assertion, fix any edge orbit $E$ and collapse
each connected component of the union of edges outside $E$.
The quotient is a nontrivial simplicial tree without inversions.
It has no fixed vertex: the preimage of one would be a proper
nonempty invariant subtree of the original tree.
Its orbits are therefore unbounded, and distance in this quotient
counts exactly the edges of $E$ in the original path $[o,go]$.
Again \eqref{found:crossing-bound} excludes a finite or two-ended
stabilizer. Apply this to the Bass--Serre tree to obtain the last
assertion.
\end{proof}

\subsection{The Hilbert space of finite-energy gradients}

We now translate the remaining contradiction assumption into the
analytic input for the compactification.
Choose a finite presentation of $\Gamma$ with relators of length
at most three; subdividing longer relators by new generators
achieves this by Tietze transformations.
Let $\mathcal C$ be its Cayley two-complex and $G$ its graph.
The vertex set is $\Gamma$. The complex $\mathcal C$ is simply
connected, $G$ is connected and locally finite, and each cellular
edge belongs to a uniformly bounded number of cell boundaries.
Keep parallel edges and loops as separate cellular edges, and
choose one orientation of each edge.

For a real function $f$ on $\Gamma$ and an oriented edge $e=uv$,
put $df(e)=f(v)-f(u)$. Define
\[
\begin{aligned}
 \mathcal D&=
 \left\{df:f:\Gamma\to\mathbb R,
                  \ \sum_{e\in E(G)}|df(e)|^2<\infty\right\},\\
 \mathcal D_0&=
 \overline{\{df:f\text{ has finite support}\}},
\end{aligned}
\]
where the closure is in $\ell^2(E(G);\mathbb R)$ and the sum
counts each cellular edge once.

\begin{lemma}\label{found:gradient}
The space $\mathcal D$ is a closed real Hilbert subspace of
$\ell^2(E(G);\mathbb R)$, and $\mathcal D_0\subseteq\mathcal D$.
If $b_1^{(2)}(\widetilde Q)>0$, then
$\mathcal D_0\ne\mathcal D$.
\end{lemma}

\begin{proof}
Write $F(\mathcal C)$ for the set of two-cells.
The cellular coboundary
\[
 d_1:\ell^2(E(G);\mathbb R)\longrightarrow
       \ell^2(F(\mathcal C);\mathbb R)
\]
is bounded: each relator has at most three edges, counted with
their cellular multiplicities, and edge--cell incidence is
uniformly bounded.
An ordinary cellular $1$-cocycle on $\mathcal C$ is a gradient:
integrate it along edge paths from a fixed vertex, using simple
connectivity to show independence of the path.
It follows that $\mathcal D=\ker d_1$, which is closed.

Likewise
$d_0:\ell^2(\Gamma;\mathbb R)\to\ell^2(E(G);\mathbb R)$
is bounded by the bounded vertex degree.
Since finite-support functions are dense in $\ell^2(\Gamma)$,
\[
 \overline{\operatorname{im}d_0}=\mathcal D_0.
\]
This does not assert that $\operatorname{im}d_0$ itself is closed.

The augmented cellular chains of $\mathcal C$ are exact in
degrees zero and one. Extend them to a free resolution by
choosing a free module $P_3$ mapping onto
$\ker(\partial_2:C_2(\mathcal C)\to C_1(\mathcal C))$, and then
mapping free modules onto successive kernels. The original cellular
terms in degrees zero, one, and two are retained and are finite free.
The singular chains of $\widetilde Q$ are free
$\mathbb Z\Gamma$-modules: a deck transformation fixing a
singular simplex fixes its image points and is therefore the
identity. Thus Lemma~\ref{found:comparison} applies.
Over $\mathbb C$, it identifies $b_1^{(2)}(\widetilde Q)$
with the dimension of
$\ker d_1\ominus\overline{\operatorname{im}d_0}$.
All matrices are real, so this complex Hilbert space is the
complexification of $\mathcal D\ominus\mathcal D_0$.
If $\mathcal D=\mathcal D_0$, its dimension would be zero.
The asserted strict inclusion follows.
\end{proof}

Henceforth assume, toward a contradiction,
$b_1^{(2)}(\widetilde Q)>0$.
We have obtained a nonzero quotient $\mathcal D/\mathcal D_0$,
together with one-endedness, the two cohomology vanishings, and
the tree-action obstruction. The next construction turns this
finite-energy information into a compact boundary.
